\begin{table*}[t]
  \centering
  \begingroup
  \small
  \setlength{\tabcolsep}{4pt}
  \renewcommand{\arraystretch}{1.16}
  \caption{Complete collected local rows for the retained branch
  \(D_3=I_3=\mathrm{ToM}\,140\).  A term \((n,e,f,d)^m\) records \(m\)
  completion factors with degree \(n=ef\), ramification index \(e\), residue
  degree \(f\), and local discriminant-exponent entry \(d\).  For each field the table
  independently gives the factor count, \(\sum mn=\CFTFieldDegree\), and
  \(\sum mfd=624\).  The \(\CFTTomOneFortyPlusDegreeOne\) degree-one factors of
  \(F_+\), absent for \(F_-\), suffice to separate the two finite
  \(\mathbf Q_3\)-algebras; these rows do not classify the individual completion
  fields.  This branch is retained alongside ToM~206, and neither branch is
  selected.}
  \label{tab:c59-local-tom140}
  \begin{tabularx}{\textwidth}{@{}c >{\centering\arraybackslash}X rrr@{}}
    \toprule
    Field & Complete \((n,e,f,d)^{\mathrm{multiplicity}}\) rows
      & Factors & \(\sum mn\) & \(\sum mfd\) \\
    \midrule
    \(F_+\) & \(\CFTTomOneFortyPlusRows\)
      & \(\CFTTomOneFortyPlusFactors\)
      & \(\CFTTomOneFortyPlusDegreeTotal\)
      & \(\CFTTomOneFortyPlusDifferentTotal\) \\
    \(F_-\) & \(\CFTTomOneFortyMinusRows\)
      & \(\CFTTomOneFortyMinusFactors\)
      & \(\CFTTomOneFortyMinusDegreeTotal\)
      & \(\CFTTomOneFortyMinusDifferentTotal\) \\
    \bottomrule
  \end{tabularx}
  \endgroup
\end{table*}
